\documentclass[11pt]{article}
\usepackage{fontspec}
\setmainfont{Times New Roman}
\setsansfont{Arial}
\setmonofont{Consolas}
\usepackage{amsmath,amssymb,mathtools}
\usepackage{geometry}
\usepackage{microtype}
\newcommand{\mo}{\mathfrak{o}}
\newcommand{\mO}{\mathfrak{O}}
\newcommand{\mT}{\mathfrak{T}}
\newcommand{\mR}{\mathfrak{R}}
\newcommand{\mq}{\mathfrak{q}}
\geometry{a4paper,margin=2.35cm}
\setlength{\parindent}{1.25em}
\setlength{\parskip}{0pt}
\makeatletter
\renewcommand\footnotesize{\@setfontsize\footnotesize{7.7}{8.7}\selectfont}
\makeatother
\emergencystretch=100em
\allowdisplaybreaks
\sloppy
\hbadness=10000
\vbadness=10000
\hfuzz=2pt
\vfuzz=10pt
\raggedbottom

\begin{document}

% Unit: Paper32 Schur paragraph v001. Direct Ukrainian translation from clean RA34 Paper32 line 8 / segment P32-S0004, checked against source scan page 1 and Claude batch OCR witness.
\setcounter{footnote}{0}

\noindent Питання про числові поля найменшого степеня над основним полем $P$, у яких незвідне над $P$ зображення скінченної групи розкладається на абсолютно незвідні складові, вперше розглянув І. Шур.\footnote{\emph{Arithmetische Untersuchungen über endliche Gruppen}, Sitzungsber. d. Berl. Ak. d. Wiss. 1906, с. 164. \emph{Beiträge zur Theorie der Gruppen linearer Substitutionen}, Transact. of the Am. Math. Soc. Ser. 2, Vol. XV, 1909, с. 159. У другій праці результати перенесено на цілком звідні гіперкомплексні системи. В один клас об'єднується зображення з усіма його перетвореними (подібними).} Нехай, наприклад, $P$ містить характер такої складової. Тоді І. Шур показав, що степінь цих полів над $P$ -- індекс -- збігається з кількістю тих абсолютно незвідних складових, які всі належать до одного й того самого класу; що, отже, те саме поле вже визначається вимогою виокремити одну абсолютно незвідну складову; далі, що степінь кожного поля над $P$, у якому відбувається такий розклад, є кратним індексу. Усі ці поля називатимемо полями розщеплення, також у випадку, коли $P$ не містить характеру. На одному прикладі І. Шур показав, що поля розщеплення найменшого степеня не обов'язково мають бути ізоморфними. Якщо $P$ не містило відповідного незвідного характеру, то поля розщеплення мають охоплювати поле такого характеру та його спряжені відносно $P$; абсолютно незвідні зображення, що належать різним спряженим характерам, хоч і спряжені, але, очевидно, належать до різних класів. Для виокремлення системи абсолютно незвідних зображень одного класу й тут досить поля відповідного характеру та одного з відповідних полів розщеплення.

\end{document}
